\documentclass[A4,svgnames,9pt,aspectratio=169]{beamer}

\usepackage[utf8]{inputenc}
\usepackage[T1]{fontenc}
\usepackage[french]{babel}
\usepackage{graphicx}
\usepackage{amsmath,amssymb}
\usepackage{booktabs}
\usepackage{tabularx}
\usepackage{xcolor}
\usepackage{tikz}
\usetikzlibrary{arrows.meta,positioning}

\usetheme{inria}
\graphicspath{{figures/}}

\hypersetup{
  colorlinks=true,
  allcolors=rouge_inria,
  pdfauthor={Louis Hauseux, Raphaël Antoine, Philippe Foucher, Pierre Charbonnier, Josiane Zerubia},
  pdftitle={CrackSAM-GeoLoRA},
  pdfsubject={Géométrie apprise dans SAM 2 : protocole, contrôle causal et résultats},
  pdfkeywords={SAM 2, CrackSAM, LoRA, Frangi, IoU tolérante, crack segmentation}
}

\definecolor{success}{HTML}{3C8D5A}
\definecolor{warning}{HTML}{D98200}
\definecolor{failure}{HTML}{C9191E}
\definecolor{inkblue}{HTML}{27348B}
\definecolor{softblue}{HTML}{E9F2F8}
\definecolor{softred}{HTML}{FBEAEC}
\definecolor{softgreen}{HTML}{EAF5EE}
\definecolor{softgray}{HTML}{F1F2F4}

\newcommand{\IoU}{\mathrm{IoU}}
\newcommand{\refmark}[1]{\textcolor{rouge_inria}{\textbf{[#1]}}}
\newcommand{\source}[1]{%
  \par\vfill{\raggedright\fontsize{5.8}{6.6}\selectfont\color{black!62}#1\par}%
}
\newcommand{\statusbadge}[2]{%
  \tikz[baseline=(n.base)]\node[rounded corners=2pt,fill=#1!13,draw=#1!55,
  inner xsep=4pt,inner ysep=1.5pt,text=#1,font=\bfseries\scriptsize] (n) {#2};%
}
\newcommand{\metriccard}[4]{%
  \begin{tikzpicture}
    \node[rounded corners=4pt,draw=#1!50,fill=#1!8,minimum width=3.05cm,
      minimum height=1.35cm,align=center,inner sep=4pt] {%
      {\scriptsize #2}\\[-1pt]
      {\Large\bfseries\color{#1} #3}\\[-1pt]
      {\fontsize{5.8}{6.5}\selectfont\color{black!62} #4}%
    };
  \end{tikzpicture}%
}
\newcommand{\imagepanel}[2]{%
  \begin{minipage}[t]{0.235\textwidth}
    \centering
    {\scriptsize\bfseries #1}\par\vspace{2pt}
    \setlength{\fboxsep}{1pt}%
    \fcolorbox{black!18}{white}{\includegraphics[width=0.95\linewidth]{#2}}
  \end{minipage}%
}
\newcommand{\compactpanel}[2]{%
  \begin{minipage}[t]{0.235\textwidth}
    \centering
    {\fontsize{6.1}{6.8}\selectfont\bfseries #1}\par\vspace{1pt}
    \setlength{\fboxsep}{0.7pt}%
    \fcolorbox{black!16}{white}{%
      \includegraphics[width=0.93\linewidth,height=1.98cm,keepaspectratio]{#2}}
  \end{minipage}%
}
% Une planche par cas : entrée, vérité terrain, référence, variante.
\newcommand{\caserow}[6]{%
  \noindent
  \begin{minipage}[c]{0.66\linewidth}
    {\scriptsize\bfseries #1}
  \end{minipage}\hfill
  \begin{minipage}[c]{0.32\linewidth}
    \raggedleft\statusbadge{#5}{#6}
  \end{minipage}\par
  \vspace{1pt}
  \compactpanel{Entrée}{#2_image.jpg}\hfill
  \compactpanel{Vérité terrain}{#2_gt.jpg}\hfill
  \compactpanel{#3}{#2_ref.jpg}\hfill
  \compactpanel{#4}{#2_var.jpg}%
}
% Quatre des onze canaux d'évidence de la même image.
\newcommand{\evidencerow}[1]{%
  \compactpanel{\texttt{frangi\_sim}}{#1_frangi.jpg}\hfill
  \compactpanel{\texttt{ofs}}{#1_ofs.jpg}\hfill
  \compactpanel{\texttt{ofa}}{#1_ofa.jpg}\hfill
  \compactpanel{\texttt{phase\_sym}}{#1_phase.jpg}%
}
% Cinq panneaux : la même image vue par la baseline, par tol3, puis par
% geo_tol3. Les deux premières colonnes viennent de la planche baseline-tol3,
% la dernière de la planche tol3-geo_tol3.
\newcommand{\widepanel}[2]{%
  \begin{minipage}[t]{0.188\textwidth}
    \centering
    {\fontsize{6.1}{6.8}\selectfont\bfseries #1}\par\vspace{1pt}
    \setlength{\fboxsep}{0.7pt}%
    \fcolorbox{black!16}{white}{%
      \includegraphics[width=0.93\linewidth,height=1.92cm,keepaspectratio]{#2}}
  \end{minipage}%
}
\newcommand{\threewayrow}[6]{%
  \noindent
  \begin{minipage}[c]{0.60\linewidth}
    {\scriptsize\bfseries #1}
  \end{minipage}\hfill
  \begin{minipage}[c]{0.38\linewidth}
    \raggedleft\statusbadge{#5}{#6}
  \end{minipage}\par
  \vspace{1pt}
  \widepanel{Entrée}{#2_image.jpg}\hfill
  \widepanel{Vérité terrain}{#2_gt.jpg}\hfill
  \widepanel{Baseline}{#2_ref.jpg}\hfill
  \widepanel{\texttt{tol3}}{#2_var.jpg}\hfill
  \widepanel{\texttt{geo\_tol3}}{#3_var.jpg}%
}
\newcommand{\synthpanel}[2]{%
  \begin{minipage}[t]{0.185\textwidth}
    \centering
    {\fontsize{6.1}{6.8}\selectfont\bfseries #1}\par\vspace{1pt}
    \setlength{\fboxsep}{0.7pt}%
    \fcolorbox{black!16}{white}{%
      \includegraphics[width=0.88\linewidth]{#2}}
  \end{minipage}%
}
\newcommand{\colourkey}{%
  \textcolor{success}{$\blacksquare$}~vrai positif \quad
  \textcolor{failure}{$\blacksquare$}~faux positif \quad
  \textcolor{inkblue}{$\blacksquare$}~manqué}

\title[CrackSAM-GeoLoRA]{CrackSAM-GeoLoRA}
\subtitle{La géométrie apprise \emph{dans} SAM 2 : protocole, contrôle causal, résultats}
\date[09/08/2026]{9 août 2026}
\author[L. \textsc{Hauseux} et al.]{%
  {\small Louis \textsc{Hauseux}, Josiane \textsc{Zerubia}}\\
  {\small Raphaël \textsc{Antoine}, Philippe \textsc{Foucher}, Pierre \textsc{Charbonnier}}%
}
\institute{Inria, équipe \textsc{Ayana} \quad---\quad Cerema, équipe \textsc{Endsum}}

\begin{document}

\frame{\titlepage}

% =============================================================================
\section{Contexte}

\begin{frame}{Quatre itérations, un même verdict}
  \centering
  \vspace{0.10cm}
  \begin{tikzpicture}[>=Latex,font=\scriptsize]
    \node[rounded corners=5pt,draw=failure!50,fill=softred,text width=2.55cm,
      minimum height=2.5cm,align=left,inner sep=6pt] at (0,0) (a) {%
      \textbf{1. Invite dense}\\[3pt]
      carte Frangi placée dans \texttt{mask\_input}\\[4pt]
      \textcolor{failure}{\textbf{$-$0,0979} d'IoU}\\
      {\fontsize{5.6}{6.2}\selectfont essai causal archivé}};
    \node[rounded corners=5pt,draw=warning!50,fill=warning!8,text width=2.55cm,
      minimum height=2.5cm,align=left,inner sep=6pt] at (3.1,0) (b) {%
      \textbf{2. Résidu sélectif}\\[3pt]
      correction révocable des logits\\[4pt]
      \textcolor{warning}{\textbf{prototype}}\\
      {\fontsize{5.6}{6.2}\selectfont robustesse aux ombres : hypothèse}};
    \node[rounded corners=5pt,draw=failure!50,fill=softred,text width=2.55cm,
      minimum height=2.5cm,align=left,inner sep=6pt] at (6.2,0) (c) {%
      \textbf{3. Arbitrage de fragments}\\[3pt]
      recoller les morceaux après coup\\[4pt]
      \textcolor{failure}{\textbf{$-$0,0066} d'IoU}\\
      {\fontsize{5.6}{6.2}\selectfont ou parité selon le pli}};
    \node[rounded corners=5pt,draw=inkblue!55,fill=softblue,text width=2.55cm,
      minimum height=2.5cm,align=left,inner sep=6pt] at (9.3,0) (d) {%
      \textbf{4. GeoLoRA}\\[3pt]
      la géométrie \textbf{apprise} dans l'encodeur\\[4pt]
      \textcolor{inkblue}{\textbf{8--9 août 2026}}\\
      {\fontsize{5.6}{6.2}\selectfont l'objet de ce rapport}};
    \draw[->,very thick,black!45] (a) -- (b);
    \draw[->,very thick,black!45] (b) -- (c);
    \draw[->,very thick,black!45] (c) -- (d);
  \end{tikzpicture}

  \vspace{0.34cm}
  \statusbadge{inkblue}{les trois premières corrigeaient un modèle gelé}
  \hspace{0.4cm}
  \statusbadge{warning}{aucune n'a produit de gain mesuré}

  \source{La carte réellement injectée en juillet était \texttt{node\_sim\_max} :
  ni MST, ni composantes, ni centralité n'intervenaient. Le résultat négatif de
  la ligne CrackSAM ne condamne donc pas le Frangi généralisé complet.}
\end{frame}

\begin{frame}{La question de cette itération}
  \begin{columns}[T,onlytextwidth]
    \begin{column}{0.475\textwidth}
      \begin{block}{Deux raisons de réessayer}
        \small
        \begin{itemize}
          \item SAM 2 ne peut structurellement pas accéder à la \textbf{portée}
            ni au \textbf{thermique} ; aucune correction en aval ne créera
            cette information.
          \item Corriger une bande de \textbf{19 px} de large demande une
            représentation \textbf{multi-échelle} qu'une retouche de logits ne
            fournit pas.
        \end{itemize}
      \end{block}
    \end{column}
    \begin{column}{0.475\textwidth}
      \begin{block}{Une raison d'en douter}
        \small
        Khánh Hà est \textbf{monomodal visible}, ses annotations sont
        \textbf{épaisses}, et la baseline est \textbf{supervisée sur ce domaine
        même}. La LoRA a peut-être déjà appris tout ce que la géométrie
        propose.
      \end{block}
    \end{column}
  \end{columns}

  \vspace{0.36cm}
  \centering
  \statusbadge{inkblue}{le protocole doit pouvoir réfuter l'hypothèse, pas seulement l'illustrer}

  \vspace{0.24cm}
  \begin{tabularx}{0.94\linewidth}{>{\bfseries}p{2.9cm}X}
    Corpus & Khánh Hà : 9\,121 train, 481 val, 1\,695 test \\
    Modèle & SAM 2 Hiera-L, LoRA q/v $r=4$, 453\,248 paramètres \\
    Question & la géométrie ajoute-t-elle une information que la LoRA n'a pas déjà~? \\
  \end{tabularx}
\end{frame}

% =============================================================================
\section{Évidence}

\begin{frame}{Ce qu'on ajoute à la Frangi--similarité : onze canaux}
  \small
  \begin{tabularx}{\linewidth}{>{\bfseries\raggedright\arraybackslash}p{2.3cm}>{\raggedright\arraybackslash}p{2.75cm}>{\raggedright\arraybackslash}X}
    \toprule
    Rôle & Canaux & Principe \\
    \midrule
    Proposer &
      \texttt{frangi\_sim} &
      le \texttt{node\_sim\_max} historique, 5 échelles --- seule carte des
      essais précédents \\
    Récuser une ombre &
      \texttt{ofs}, \texttt{ofa}, \texttt{abs\_energy} &
      Oriented Flux Symmetry : tenseur de flux sur un disque, et antisymétrie
      sur l'anneau comme confiance \refmark{5} \\
    Trancher ligne \emph{vs} marche &
      \texttt{even\_odd}, \texttt{dbic}, \texttt{profile} &
      réponse paire moins impaire ; sélection de modèle $\Delta$BIC entre
      $H_0$ et $H_1$ ; deux flancs plus clairs au même rayon \\
    Polarité sombre &
      \texttt{phase\_sym} &
      Gabor en quadrature ; une marche est rejetée par $|\text{odd}|$
      \refmark{6} \\
    Géométrie explicite &
      \texttt{cos2}$\theta$, \texttt{sin2}$\theta$, \texttt{scale} &
      orientation et rayon gagnant, bornés dans $[0,1]$ \\
    \bottomrule
  \end{tabularx}

  \vspace{0.28cm}
  \centering
  \statusbadge{failure}{jamais multipliés entre eux}
  \hspace{0.35cm}
  \statusbadge{inkblue}{séparés jusqu'à l'encodeur}

  \source{L'étude filtre-seul du 8 août montre pourquoi : la fusion
  \emph{multiplicative} anti-ombre $\times$ Frangi retombe à environ
  \textbf{10~\%} de rétention quand la fissure passe sous l'ombre, contre
  \textbf{89,5~\%} pour OFS seul. Sur 24 paires ligne/marche, la fuite médiane
  d'OFS vaut \textbf{0,133} contre \textbf{1,000} pour la carte historique.}
\end{frame}

\begin{frame}{Les onze canaux sur une image réelle}
  \centering
  \vspace{0.06cm}
  \imagepanel{Image}{b2t_gain_ct6774_image.jpg}\hfill
  \imagepanel{\texttt{frangi\_sim}}{b2t_gain_ct6774_frangi.jpg}\hfill
  \imagepanel{\texttt{ofs}}{b2t_gain_ct6774_ofs.jpg}\hfill
  \imagepanel{\texttt{phase\_sym}}{b2t_gain_ct6774_phase.jpg}

  \vspace{0.34cm}
  \begin{columns}[T,onlytextwidth]
    \begin{column}{0.31\textwidth}
      \centering\statusbadge{failure}{bande épaisse, hors sujet}\\[2pt]
      {\fontsize{6.4}{7.4}\selectfont \texttt{frangi\_sim} est accordé sur
      19,1 px : il fusionne le réseau en une seule coulée.}
    \end{column}
    \begin{column}{0.31\textwidth}
      \centering\statusbadge{warning}{presque vide}\\[2pt]
      {\fontsize{6.4}{7.4}\selectfont \texttt{ofs} s'abstient --- ce qui est le
      comportement voulu, mais il ne propose alors plus rien.}
    \end{column}
    \begin{column}{0.31\textwidth}
      \centering\statusbadge{warning}{saturé d'artefacts}\\[2pt]
      {\fontsize{6.4}{7.4}\selectfont \texttt{phase\_sym} répond aux étoiles de
      Gabor autant qu'aux fissures.}
    \end{column}
  \end{columns}

  \source{\texttt{cracktree200\_6774}, vérité terrain à 0,8~\% de pixels. Les
  fissures y font 1 à 2 px : c'est le régime que le réaccordage sur la largeur
  \emph{moyenne} a sur-corrigé.}
\end{frame}

\begin{frame}{Des échelles mesurées, pas héritées}
  \begin{columns}[T,onlytextwidth]
    \begin{column}{0.47\textwidth}
      \small
      \begin{tabularx}{\linewidth}{>{\bfseries}p{1.9cm}>{\raggedright\arraybackslash}X}
        \toprule
        Filtre & Paramètre \\
        \midrule
        Frangi & $\sigma=\{1{,}5\,;3\,;5\,;8\,;12\}$ \\
        OFS & rayons $\{2,3,4,6,8\}$ \\
        Phase & $\lambda=\{5\,;8\,;12\,;18\}$ \\
        Profil, $\Delta$BIC & $\sigma=\{1{,}2\,;2\,;3\,;4{,}5\,;6\}$ \\
        \bottomrule
      \end{tabularx}

      \vspace{0.2cm}
      {\fontsize{6.6}{7.6}\selectfont L'évidence est calculée à
      \textbf{224 px} : une fissure de 19,1 px à 448 y fait 9,6 px, au centre
      de la plage validée. À 448, couvrir la même structure demanderait des
      noyaux de Gabor de $143\times143$.}
    \end{column}
    \begin{column}{0.47\textwidth}
      \begin{block}{L'erreur que la moyenne a masquée}
        \small
        La largeur de \textbf{19,1 px} est une médiane de corpus. Les scènes
        \texttt{cracktree200}, \texttt{GAPS384} et \texttt{Rissbilder} ont des
        vérités terrain à \textbf{0,4--0,8~\%} de pixels : des fissures
        \textbf{fines}. En réaccordant sur 19,1 px, on a sur-corrigé pour une
        partie du corpus.
      \end{block}
      \centering
      \statusbadge{warning}{une évidence multi-échelle, et non mono-largeur}
    \end{column}
  \end{columns}

  \source{Échelles gelées dans \texttt{manifest\_train.json} :
  \texttt{target\_width\_px} $=9{,}55$, 11 canaux, 9\,605 caches d'évidence
  calculés (environ 19 s par image, sur CPU).}
\end{frame}

% =============================================================================
\section{Intégration}

\begin{frame}{Comment la géométrie entre dans SAM 2}
  \centering
  \vspace{0.02cm}
  \begin{tikzpicture}[>=Latex,font=\scriptsize]
    \node[rounded corners=4pt,draw=black!30,fill=softgray,text width=2.0cm,
      align=center,inner sep=4pt] at (0,0.95) (rgb) {Image RGB\\$448^2$};
    \node[rounded corners=4pt,draw=warning!55,fill=warning!8,text width=2.35cm,
      align=center,inner sep=4pt] at (3.0,0.95) (hiera) {%
      \textbf{SAM 2 Hiera-L}\\LoRA q/v $r=4$\\453\,248 params};
    \node[rounded corners=4pt,draw=black!30,fill=softgray,text width=2.0cm,
      align=center,inner sep=4pt] at (0,-1.15) (geo) {11 canaux\\$224^2$};
    \node[rounded corners=4pt,draw=inkblue!55,fill=softblue,text width=2.35cm,
      align=center,inner sep=4pt] at (3.0,-1.15) (enc) {%
      \textbf{Encodeur géométrique}\\290\,801 params};
    \node[rounded corners=4pt,draw=inkblue!45,fill=white,text width=3.35cm,
      align=left,inner sep=4pt] at (7.1,-1.15) (inj) {%
      \textcolor{inkblue}{\textbf{3 projections init. à ZÉRO}}\\[1.5pt]
      $+$ \texttt{high\_res[0]} \hfill $32{\times}256^2$\\
      $+$ \texttt{high\_res[1]} \hfill $64{\times}128^2$\\
      $+$ \texttt{embeddings} \hfill $256{\times}64^2$};
    \node[rounded corners=4pt,draw=success!50,fill=softgreen,text width=2.35cm,
      align=center,inner sep=4pt] at (10.7,0.95) (dec) {%
      \textbf{Mask decoder}\\logits $448^2$};

    \draw[->,thick] (rgb) -- (hiera);
    \draw[->,thick] (geo) -- (enc);
    \draw[->,thick] (enc) -- (inj);
    \draw[->,thick] (hiera) -- (dec);
    \draw[->,thick] (inj.east) -| (dec.south);
  \end{tikzpicture}

  \vspace{0.30cm}
  \statusbadge{inkblue}{à l'initialisation, le modèle \emph{est} la baseline gelée}
  \hspace{0.3cm}
  \statusbadge{success}{\texttt{evidence=None} reste bit à bit}
  \hspace{0.3cm}
  \statusbadge{failure}{\texttt{mask\_input} reste un contrôle négatif}

  \source{Les corrections sont \emph{additives} et injectées à trois
  résolutions, parce que les corridors mono-échelle des essais précédents ne
  couvraient que 1,8~\% de la vérité terrain contre 5,7~\% attendus.}
\end{frame}

\begin{frame}{Le point selle qui a coûté une heure de GPU}
  \centering
  \vspace{0.10cm}
  \begin{tikzpicture}[>=Latex,font=\small]
    \node[rounded corners=5pt,draw=failure!50,fill=softred,text width=4.6cm,
      minimum height=2.35cm,align=center,inner sep=6pt] at (0,0) {%
      \textbf{Version fautive}\\[4pt]
      $\mathrm{d}z = \gamma \cdot \mathrm{proj}(x)$\\[4pt]
      $\gamma = 0$ \textbf{et} $\mathrm{proj} = 0$\\[3pt]
      {\scriptsize les \emph{deux} gradients s'annulent}};
    \node[rounded corners=5pt,draw=success!50,fill=softgreen,text width=4.6cm,
      minimum height=2.35cm,align=center,inner sep=6pt] at (6.4,0) {%
      \textbf{Version corrigée}\\[4pt]
      $\gamma = 1$, $\mathrm{proj} = 0$\\[4pt]
      $\mathrm{d}z = 0$ à l'initialisation\\[3pt]
      {\scriptsize et le gradient existe}};
    \draw[->,very thick,black!45] (2.35,0) -- (4.05,0);
  \end{tikzpicture}

  \vspace{0.34cm}
  \begin{columns}[T,onlytextwidth]
    \begin{column}{0.47\textwidth}
      \centering
      \statusbadge{failure}{constaté en réel}\\[3pt]
      {\scriptsize activation $=0{,}0000$ après une époque complète, et la
      variante géométrique \emph{numériquement identique} à sa version sans
      géométrie.}
    \end{column}
    \begin{column}{0.47\textwidth}
      \centering
      \statusbadge{success}{test de régression}\\[3pt]
      {\scriptsize\texttt{test\_adapter\_gradients\_are\_not\_both\_dead\_}
      \texttt{at\_initialisation} --- l'un des 15 tests du module.}
    \end{column}
  \end{columns}

  \source{Trois autres incidents ont coûté du GPU sans changer les
  conclusions : un correctif transféré mais jamais extrait, deux entraînements
  concurrents sur le même GPU, et un pilote NVIDIA cassé par un changement de
  noyau. Le seul qui aurait pu \emph{fausser} un résultat est le point selle.}
\end{frame}

% =============================================================================
\section{Protocole}

\begin{frame}{Juger honnêtement : six exigences}
  \small
  \begin{tabularx}{\linewidth}{>{\bfseries\raggedright\arraybackslash}p{3.45cm}>{\raggedright\arraybackslash}X}
    \toprule
    Budget strictement égal &
      5 époques, graine 3407, même ordonnancement, même départ depuis la LoRA
      convergée. La \textbf{baseline est réentraînée} dans ces conditions :
      l'archivée valait 70 époques, elle n'est pas comparable. \\
    Un facteur par barreau &
      \texttt{cldice} isole la perte de continuité ; \texttt{geo} y ajoute la
      géométrie ; \texttt{tol3} isole la tolérance \emph{sans} clDice. \\
    Un contrôle apparié &
      \texttt{geo\_tol3\_permuted} : même capacité, même perte, même graine ;
      seul l'\textbf{alignement image~$\leftrightarrow$~évidence} est détruit.
      Sans lui, on ne sépare pas \og la géométrie aide\fg{} de \og les paramètres
      aident\fg{}. \\
    Un test de nécessité d'entrée &
      à l'évaluation, \texttt{evidence=None} restitue la voie baseline bit à
      bit : si les scores ne bougent pas, l'adapter est décoratif. \\
    Des deltas appariés par image &
      gains \emph{et} pertes comptés séparément : une moyenne ne dit pas que
      866 images sur 1\,695 ont perdu. \\
    La métrique choisie \emph{avant} &
      deux conventions de tolérance ne classent pas pareil ; choisir après
      avoir vu les chiffres, c'est choisir son résultat. \\
    \bottomrule
  \end{tabularx}

  \vspace{0.24cm}
  \centering
  \statusbadge{warning}{ce qui manque : IC95 par bootstrap groupé, \texttt{tol5}, et le contrôle permuté de la famille clDice}

  \source{Exécution des 8--9 août 2026 sur RTX PRO 6000 Blackwell (97,9 Go).
  Chaque époque écrit un état complet, optimiseur compris : la reprise après
  préemption Spot repart à l'époque suivante.}
\end{frame}

\begin{frame}{L'échelle d'ablations}
  \small
  \begin{tabularx}{\linewidth}{c>{\ttfamily}p{2.75cm}p{2.3cm}>{\raggedright\arraybackslash}p{1.7cm}r>{\raggedright\arraybackslash}X}
    \toprule
    \# & \normalfont Variante & Perte & Évidence & Params &
      \normalfont Ce qu'elle isole \\
    \midrule
    0 & baseline & CrackSAM & --- & 453\,248 & l'ancre à budget égal \\
    1 & cldice & $+$ soft-clDice & --- & 453\,248 & la continuité seule \\
    2 & geo & $+$ soft-clDice & alignée & 744\,049 & la géométrie \\
    3 & tol3 & $+$ tolérante $k{=}3$ & --- & 453\,248 & la tolérance seule \\
    4 & geo\_tol3 & $+$ tolérante $k{=}3$ & alignée & 744\,049 &
      la géométrie, sans \mbox{clDice} \\
    5 & geo\_tol3\_permuted & $+$ tolérante $k{=}3$ &
      \textbf{d'une autre image} & 744\,049 & la capacité seule \\
    \bottomrule
  \end{tabularx}

  \vspace{0.30cm}
  \centering
  \metriccard{inkblue}{Durée par variante}{40--41 min}{61 min pour \texttt{geo}}
  \hspace{0.25cm}
  \metriccard{success}{Coût de l'adapter}{$+1$ min}{744\,049 contre 453\,248 params}
  \hspace{0.25cm}
  \metriccard{warning}{Époques}{5}{et non 20, l'optimum archivé}

  \source{Les 21 minutes supplémentaires de \texttt{geo} viennent de
  \texttt{soft-clDice}, dont la squelettisation douce coûte dix passes de
  \emph{pooling} par échantillon --- et non de la géométrie.}
\end{frame}

% =============================================================================
\section{Mesurer}

\begin{frame}{L'IoU stricte mesure surtout un biais de largeur}
  \begin{columns}[T,onlytextwidth]
    \begin{column}{0.45\textwidth}
      \centering
      {\small\bfseries La vérité terrain contre elle-même}\\[3pt]
      \small
      \begin{tabularx}{\linewidth}{Xr}
        \toprule
        Perturbation du GT & IoU \\
        \midrule
        dilaté de 1 px & 0,881 \\
        érodé de 1 px & 0,843 \\
        dilaté de 2 px & 0,799 \\
        \bottomrule
      \end{tabularx}
      \\[5pt]
      {\fontsize{6.4}{7.4}\selectfont Un décalage d'un pixel coûte
      \textbf{douze points} d'IoU --- davantage que tous les écarts entre
      variantes mesurés ici.}
    \end{column}
    \begin{column}{0.51\textwidth}
      \centering
      {\small\bfseries Baseline : précision et rappel tolérants}\\[3pt]
      \small
      \begin{tabularx}{\linewidth}{Xrrrr}
        \toprule
        $k$ & 0 & 1 & 2 & 3 \\
        \midrule
        Précision & 0,759 & \textbf{0,888} & 0,923 & 0,941 \\
        Rappel & 0,761 & 0,812 & 0,848 & 0,879 \\
        \bottomrule
      \end{tabularx}
      \\[5pt]
      {\fontsize{6.4}{7.4}\selectfont Un seul pixel de tolérance fait gagner
      \textbf{$+0{,}129$} à la précision et seulement $+0{,}051$ au rappel :
      l'essentiel de ce que l'IoU stricte compte comme faux positifs est du
      \textbf{débordement de frontière}, pas de la fausse détection.}
    \end{column}
  \end{columns}

  \vspace{0.30cm}
  \centering
  \statusbadge{failure}{comparer deux méthodes à $\pm0{,}02$ d'IoU stricte revient largement à comparer leurs largeurs}
\end{frame}

\begin{frame}{Deux conventions de tolérance, deux classements}
  \begin{columns}[T,onlytextwidth]
    \begin{column}{0.475\textwidth}
      \begin{block}{\texttt{buffered} --- à privilégier}
        \small
        \[
          P_k=\frac{|P\cap \delta_k G|}{|P|},\qquad
          R_k=\frac{|G\cap \delta_k P|}{|G|}
        \]
        \[
          \IoU_k = \frac{F1_k}{2-F1_k}
        \]
        {\fontsize{6.4}{7.4}\selectfont Ne dilate \emph{jamais} les deux
        masques : mesure une distance d'appariement. Dilatation euclidienne,
        donc isotrope.}
      \end{block}
    \end{column}
    \begin{column}{0.475\textwidth}
      \begin{block}{\texttt{dilate\_both} --- convention EUVIP du dépôt}
        \small
        \[
          \IoU_k = \IoU\bigl(\delta_k P,\ \delta_k G\bigr)
        \]
        \vspace{-0.15cm}
        {\fontsize{6.4}{7.4}\selectfont C'est \texttt{thicken(sk, 6)} puis
        Jaccard. Elle épaissit les deux masques, donc elle favorise les
        prédictions déjà larges, et sature vers 1.}
      \end{block}
    \end{column}
  \end{columns}

  \vspace{0.24cm}
  \centering
  \small
  \begin{tabularx}{0.90\linewidth}{Xrrrrrr}
    \toprule
    $k=8$ & \ttfamily baseline & \ttfamily cldice & \ttfamily geo &
      \ttfamily tol3 & \ttfamily geo\_tol3 & \ttfamily perm. \\
    \midrule
    \texttt{buffered} & 0,9213 & 0,9341 & 0,9347 & 0,9346 & 0,9374 &
      \textbf{0,9378} \\
    \texttt{dilate\_both} & 0,7720 & 0,8175 & \textbf{0,8181} & 0,7943 &
      0,7960 & 0,7974 \\
    \bottomrule
  \end{tabularx}

  \vspace{0.18cm}
  \statusbadge{failure}{la même mesure à la même tolérance donne deux vainqueurs différents}
\end{frame}

\begin{frame}{Ce que la tolérance pardonne --- et ce qu'elle ne pardonne pas}
  \centering
  \vspace{0.02cm}
  \synthpanel{parfait}{synth_perfect.png}\hfill
  \synthpanel{décalé 2 px}{synth_shifted.png}\hfill
  \synthpanel{2$\times$ trop large}{synth_wide.png}\hfill
  \synthpanel{rompu 4 px}{synth_break4.png}\hfill
  \synthpanel{rompu 20 px}{synth_break20.png}

  \vspace{0.20cm}
  \small
  \begin{tabularx}{0.88\linewidth}{Xrrrrrr}
    \toprule
    IoU tampon & $k{=}0$ & $k{=}1$ & $k{=}2$ & $k{=}3$ & $k{=}5$ & $k{=}8$ \\
    \midrule
    prédiction parfaite & 1,000 & 1,000 & 1,000 & 1,000 & 1,000 & 1,000 \\
    décalée de 2 px & 0,200 & 0,500 & \textbf{1,000} & 1,000 & 1,000 & 1,000 \\
    deux fois trop large & 0,429 & 0,714 & \textbf{1,000} & 1,000 & 1,000 & 1,000 \\
    rompue sur 4 px & 0,917 & 0,958 & \textbf{1,000} & 1,000 & 1,000 & 1,000 \\
    \textbf{rompue sur 20 px} & 0,583 & 0,625 & 0,667 & 0,708 & 0,792 &
      \textbf{0,917} \\
    \bottomrule
  \end{tabularx}

  \vspace{0.20cm}
  \statusbadge{failure}{l'IoU stricte sanctionne un décalage de 2 px (0,200) bien plus qu'une rupture de 20 px (0,583)}

  \source{Ligne verticale de 3 px sur 48, dans une image de $64^2$ ; scores
  recalculés par \texttt{scripts/05\_tolerant\_iou.py} lui-même. La tolérance
  est une distance d'appariement : elle efface donc aussi une rupture, mais
  seulement tant que celle-ci est plus courte que $2k$. Une coupure franche
  reste pénalisée à toutes les tolérances mesurées.}
\end{frame}

\begin{frame}{La perte tolérante, contrepartie différentiable}
  \begin{columns}[T,onlytextwidth]
    \begin{column}{0.50\textwidth}
      \vspace{0.06cm}
      \begin{tikzpicture}
        \node[rounded corners=4pt,draw=black!25,fill=softgray,align=left,
          inner sep=6pt,text width=5.6cm,font=\fontsize{7.4}{9.4}\selectfont] {%
          $G_k = \texttt{soft\_dilate}(G,\,k)$ \hfill {\tiny max-pool}\\
          $P_k = \texttt{soft\_dilate}(P,\,k)$\\
          $\text{prec} = \sum (P \cdot G_k)\,/\,\sum P$\\
          $\text{rec}\ \, = \sum (G \cdot P_k)\,/\,\sum G$\\
          $L = 1 - 2\,\text{prec}\cdot\text{rec}\,/\,(\text{prec}+\text{rec})$};
      \end{tikzpicture}
      \\[3pt]
      {\fontsize{6.4}{7.4}\selectfont Même machinerie morphologique que
      \texttt{soft\_skeleton}. La dilatation est \emph{carrée} et non
      euclidienne : le rayon n'est pas une grandeur à estimer, c'est une
      tolérance choisie.}
    \end{column}
    \begin{column}{0.46\textwidth}
      \centering
      \small
      \setlength{\tabcolsep}{3pt}%
      \begin{tabularx}{\linewidth}{>{\raggedright\arraybackslash}p{2.1cm}rrr}
        \toprule
        $1-F1$ tolérant & $k{=}0$ & $k{=}1$ & $k{=}3$ \\
        \midrule
        parfait & 0,0000 & 0,0000 & 0,0000 \\
        décalé de 2 px & 0,6621 & 0,3310 & \textbf{0,0000} \\
        2$\times$ trop large & 0,3983 & 0,1661 & \textbf{0,0000} \\
        rompu, 20 px & 0,2609 & 0,2288 & \textbf{0,1694} \\
        \bottomrule
      \end{tabularx}
      \\[6pt]
      \statusbadge{success}{le placement est pardonné, la rupture non}
    \end{column}
  \end{columns}

  \vspace{0.26cm}
  \centering
  \statusbadge{inkblue}{les variantes \texttt{tol*} n'incluent volontairement \emph{pas} clDice : on isole la tolérance seule}

  \source{Trois variantes ont été câblées dans le CLI d'entraînement :
  \texttt{tol3}, \texttt{tol5} et \texttt{geo\_tol3}. \texttt{tol5} a été
  sacrifié pour faire place au contrôle causal, qui importait davantage : la
  sensibilité au rayon reste donc inconnue.}
\end{frame}

% =============================================================================
\section{Résultats}

\begin{frame}{Les six barreaux, en IoU stricte}
  \small
  \begin{tabularx}{\linewidth}{>{\ttfamily}p{2.85cm}Xrrrrrr}
    \toprule
    \normalfont Variante & & IoU & Précision & Rappel & Couv. sq. &
      Comp. & Sans évid. \\
    \midrule
    baseline & & 0,6241 & 0,7588 & 0,7607 & 0,6792 & 2,76 & --- \\
    cldice & & 0,6066 & 0,6738 & 0,8567 & \textbf{0,8254} & 2,65 & --- \\
    geo & & 0,6083 & 0,6742 & \textbf{0,8590} & \textbf{0,8263} & 2,52 & 0,6138 \\
    \textbf{tol3} & & \textbf{0,6276} & 0,7400 & 0,7940 & 0,7238 & 3,31 & --- \\
    geo\_tol3 & & 0,6270 & 0,7432 & 0,7869 & 0,7162 & 5,97 & 0,6281 \\
    geo\_tol3\_permuted & & 0,6265 & 0,7383 & 0,7925 & 0,7232 & 5,66 & 0,6281 \\
    \bottomrule
  \end{tabularx}

  \vspace{0.26cm}
  \centering
  \metriccard{inkblue}{Baseline à budget égal}{0,6241}{ancre, 5 époques}
  \hspace{0.25cm}
  \metriccard{success}{\texttt{tol3}}{0,6276}{$+0{,}0035$, sans géométrie}
  \hspace{0.25cm}
  \metriccard{failure}{Contrôle permuté}{0,6265}{à $0{,}0005$ de \texttt{geo\_tol3}}

  \source{Vérité terrain : 57,5 composantes connexes par image en moyenne,
  contre 2,5 à 6,0 pour les prédictions --- toutes les variantes
  sous-segmentent massivement en \emph{nombre} de composantes, y compris
  celles qui optimisent la continuité.}
\end{frame}

\begin{frame}{IoU tampon, six tolérances : le classement s'inverse}
  \centering
  \small
  \begin{tabularx}{0.96\linewidth}{Xrrrrrr}
    \toprule
    $k$ & \ttfamily baseline & \ttfamily cldice & \ttfamily geo &
      \ttfamily tol3 & \ttfamily geo\_tol3 & \ttfamily geo\_tol3\_perm. \\
    \midrule
    \textbf{0} & 0,6241 & 0,6066 & 0,6083 & \textbf{0,6276} & 0,6270 & 0,6265 \\
    1 & 0,7407 & 0,7574 & 0,7593 & 0,7586 & 0,7597 & \textbf{0,7601} \\
    2 & 0,7971 & 0,8213 & \textbf{0,8231} & 0,8186 & 0,8206 & 0,8214 \\
    \textbf{3} & 0,8396 & 0,8659 & \textbf{0,8674} & 0,8607 & 0,8635 & 0,8644 \\
    5 & 0,8901 & 0,9110 & \textbf{0,9119} & 0,9072 & 0,9101 & 0,9108 \\
    8 & 0,9213 & 0,9341 & 0,9347 & 0,9346 & 0,9374 & \textbf{0,9378} \\
    \bottomrule
  \end{tabularx}

  \vspace{0.28cm}
  \begin{columns}[T,onlytextwidth]
    \begin{column}{0.31\textwidth}
      \centering\statusbadge{success}{à $k=0$, \texttt{tol3} gagne}\\[2pt]
      {\fontsize{6.4}{7.4}\selectfont Cesser de facturer une frontière que
      l'annotation ne définit pas \textbf{libère de la capacité utile}.}
    \end{column}
    \begin{column}{0.31\textwidth}
      \centering\statusbadge{inkblue}{dès $k=1$, tout dépasse la baseline}\\[2pt]
      {\fontsize{6.4}{7.4}\selectfont Le verdict \og aucune variante ne bat la
      baseline\fg{} ne tenait qu'à tolérance nulle. À $k=3$ : $+0{,}0278$ pour
      \texttt{geo}.}
    \end{column}
    \begin{column}{0.31\textwidth}
      \centering\statusbadge{failure}{mais le contrôle mène 5 fois sur 6}\\[2pt]
      {\fontsize{6.4}{7.4}\selectfont \texttt{geo\_tol3\_permuted} est devant
      \texttt{geo\_tol3} à toutes les tolérances sauf $k=0$.}
    \end{column}
  \end{columns}
\end{frame}

\begin{frame}{Le contrôle qui tranche}
  \begin{columns}[T,onlytextwidth]
    \begin{column}{0.46\textwidth}
      \centering
      \small
      \setlength{\tabcolsep}{3pt}%
      \begin{tabularx}{\linewidth}{Xrr}
        \toprule
        $k$ &
          \shortstack[r]{\ttfamily tol3\\[-1pt] $-$ \ttfamily baseline} &
          \shortstack[r]{alignée\\[-1pt] $-$ permutée} \\
        \midrule
        0 & $+0{,}0035$ & $+0{,}0004$ \\
        1 & $+0{,}0179$ & $-0{,}0005$ \\
        2 & $+0{,}0215$ & $-0{,}0008$ \\
        3 & $+0{,}0212$ & $-0{,}0009$ \\
        5 & $+0{,}0171$ & $-0{,}0007$ \\
        8 & $+0{,}0133$ & $-0{,}0004$ \\
        \bottomrule
      \end{tabularx}
      \\[5pt]
      {\fontsize{6.4}{7.4}\selectfont L'effet de la perte est \textbf{vingt à
      quarante fois} plus grand que l'écart entre géométrie alignée et
      géométrie d'une autre image.}
    \end{column}
    \begin{column}{0.50\textwidth}
      \begin{block}{Ce que le contrôle établit}
        \small
        L'adapter \textbf{s'active} : ses projections croissent jusqu'à
        $1{,}79\times10^{-3}$, et le test de nécessité d'entrée est positif
        (\texttt{evidence=None} change le score). Mais ce qu'il apprend
        \textbf{ne dépend pas} de la correspondance entre l'évidence et
        l'image.
      \end{block}
      \centering
      \metriccard{failure}{Deltas appariés, par image}{$-0{,}0007$}{603 gains,
      866 pertes, 226 nuls}
    \end{column}
  \end{columns}

  \vspace{0.22cm}
  \centering
  \statusbadge{failure}{aligner la géométrie sur l'image ne sert à rien}
  \hspace{0.3cm}
  \statusbadge{inkblue}{le déficit de \texttt{geo\_tol3} vient des 290\,801 paramètres, pas d'un mauvais usage}
\end{frame}

% =============================================================================
\section{Cas}

\begin{frame}{Ce que la perte tolérante change, et ce qu'elle coûte}
  \caserow{\texttt{cracktree200\_6774} --- la baseline manque
    \emph{intégralement} le réseau}{b2t_gain_ct6774}{Baseline}{Avec
    \texttt{tol3}}{success}{0,000 $\rightarrow$ 0,219}

  \vspace{0.26cm}
  \caserow{\texttt{CRACK500\_20160405\_171336} --- annotation épaisse
    (10,6~\% des pixels)}{b2t_loss_crack171336}{Baseline}{Avec
    \texttt{tol3}}{failure}{0,875 $\rightarrow$ 0,510}

  \vspace{0.16cm}
  \centering{\fontsize{6.4}{7.4}\selectfont \colourkey}

  \source{Les gains se concentrent sur \texttt{cracktree200}, fissures fines et
  peu contrastées : le régime où exiger une frontière au pixel près est le plus
  coûteux et le moins pertinent. Les pertes se concentrent sur les annotations
  larges, où la tolérance autorise un débordement que l'IoU stricte sanctionne.}
\end{frame}

\begin{frame}{La famille clDice : plus de rappel, moins de précision}
  \caserow{\texttt{cracktree200\_6774} --- le réseau fin est
    retrouvé}{b2g_gain_ct6774}{Baseline}{Avec \texttt{geo}}{success}{0,000
    $\rightarrow$ 0,228}

  \vspace{0.26cm}
  \caserow{\texttt{CFD\_080} --- le débordement autour d'une fissure déjà bien
    détectée}{b2g_loss_cfd080}{Baseline}{Avec
    \texttt{geo}}{failure}{0,523 $\rightarrow$ 0,343}

  \vspace{0.16cm}
  \centering
  \statusbadge{success}{couverture du squelette : 0,679 $\rightarrow$ 0,826}
  \hspace{0.3cm}
  \statusbadge{failure}{précision : 0,759 $\rightarrow$ 0,674}

  \source{C'est le même mécanisme dans les deux sens : \texttt{soft-clDice}
  pousse à couvrir les squelettes, donc récupère les réseaux manqués et
  déborde là où la baseline était déjà juste. Ce que l'IoU stricte comptait
  comme faux positifs est en grande partie ce halo rouge.}
\end{frame}

\begin{frame}{Ce que la géométrie change}
  \caserow{\texttt{Rissbilder\_for\_Florian} --- le plus grand gain de la
    paire}{t2g_gain_riss}{\texttt{tol3}}{\texttt{geo\_tol3}}{success}{0,153
    $\rightarrow$ 0,405}

  \vspace{0.26cm}
  \caserow{\texttt{GAPS384\_0608} --- la contrepartie, aussi
    fréquente}{t2g_loss_gaps608}{\texttt{tol3}}{\texttt{geo\_tol3}}{failure}{0,443
    $\rightarrow$ 0,341}

  \vspace{0.16cm}
  \centering
  \statusbadge{inkblue}{603 gains contre 866 pertes : les deux colonnes existent, en nombre comme en amplitude}

  \source{Aucun de ces deux cas n'est atypique : la distribution des deltas
  appariés est un pic étroit centré sur zéro, avec une queue de part et
  d'autre.}
\end{frame}

\begin{frame}{Les deux plus grands gains, vus de bout en bout}
  \threewayrow{\texttt{CRACK500\_20160405\_171336}}{b2t_loss_crack171336}{t2g_gain_crack171336}{}{warning}{0,875 $\rightarrow$ 0,510 $\rightarrow$ 0,807}

  \vspace{0.24cm}
  \threewayrow{\texttt{Rissbilder\_for\_Florian}}{b2t_loss_riss}{t2g_gain_riss}{}{warning}{0,422 $\rightarrow$ 0,153 $\rightarrow$ 0,405}

  \vspace{0.14cm}
  \centering
  \statusbadge{failure}{la géométrie récupère ce que la tolérance avait coûté, et s'arrête juste avant la baseline}

  \source{Ces deux images sont exactement les deux plus grandes \emph{pertes}
  de \texttt{tol3} face à la baseline ($-0{,}364$ et $-0{,}269$) et les deux
  plus grands \emph{gains} de \texttt{geo\_tol3} face à \texttt{tol3}
  ($+0{,}296$ et $+0{,}252$). Une galerie de gains construite sur ces deux
  images raconterait donc l'inverse de ce que dit le corpus.}
\end{frame}

\begin{frame}{Le piège de la galerie sélective}
  \centering
  \vspace{0.04cm}
  \evidencerow{t2g_gain_riss}

  \vspace{0.22cm}
  \small
  \begin{tabularx}{0.90\linewidth}{Xrrrr}
    \toprule
    Les deux plus grands gains de \texttt{geo\_tol3} & \ttfamily baseline &
      \ttfamily tol3 & \ttfamily geo\_tol3 & $\Delta$ \\
    \midrule
    \texttt{CRACK500\_20160405\_171336} & \textbf{0,875} & 0,510 & 0,807 &
      $+0{,}296$ \\
    \texttt{Rissbilder\_for\_Florian} & \textbf{0,422} & 0,153 & 0,405 &
      $+0{,}252$ \\
    \bottomrule
  \end{tabularx}

  \vspace{0.22cm}
  \begin{columns}[T,onlytextwidth]
    \begin{column}{0.47\textwidth}
      \centering\statusbadge{failure}{ce sont exactement les deux plus grandes
      \emph{pertes} de \texttt{tol3}}\\[3pt]
      {\fontsize{6.4}{7.4}\selectfont La géométrie ne fait que rendre ce que la
      perte tolérante avait coûté --- et reste \textbf{sous la baseline} dans
      les deux cas.}
    \end{column}
    \begin{column}{0.47\textwidth}
      \centering\statusbadge{warning}{et aucun canal ne désigne la
      fissure}\\[3pt]
      {\fontsize{6.4}{7.4}\selectfont Sur \texttt{Rissbilder} :
      \texttt{frangi\_sim} est une tache sur le bord supérieur, \texttt{ofa}
      sature de coulures, \texttt{phase\_sym} n'est qu'un entrelacs.}
    \end{column}
  \end{columns}

  \source{Le contrôle permuté obtient la \emph{même} performance moyenne : il
  s'agit de variance d'entraînement, pas d'un effet de la géométrie. Une figure
  ne remplace pas un contrôle.}
\end{frame}

% =============================================================================
\section{Bilan}

\begin{frame}{Trois conclusions fausses, corrigées par la mesure}
  \small
  \begin{tabularx}{\linewidth}{>{\bfseries\raggedright\arraybackslash}p{3.5cm}>{\raggedright\arraybackslash}X}
    \toprule
    \og L'adapter coûte 2,5$\times$ plus cher\fg{} &
      Faux : c'était la contention entre deux entraînements concurrents sur le
      même GPU. Mesuré proprement, \texttt{geo\_tol3} prend 41 min contre 40
      pour \texttt{tol3}. Les 61 min de \texttt{geo} viennent de
      \texttt{soft-clDice}. \\
    \og \texttt{soft-clDice} dégrade\fg{} &
      Faux sous une métrique adaptée : $-0{,}0175$ en IoU stricte,
      $+0{,}0263$ à $k=3$. Ce que l'IoU stricte comptait comme faux positifs
      était du débordement de frontière. \\
    \og Aucune variante ne bat la baseline\fg{} &
      Vrai à tolérance nulle \emph{uniquement}. \texttt{tol3} la bat même en
      strict ($+0{,}0035$), et toute la famille tolérante la dépasse dès
      $k=1$. \\
    \bottomrule
  \end{tabularx}

  \vspace{0.30cm}
  \centering
  \statusbadge{failure}{trois fois, une conclusion tirée trop tôt d'une mesure inadaptée}

  \vspace{0.16cm}
  \statusbadge{success}{le seul garde-fou qui a fonctionné : mesurer autrement, et exécuter le contrôle}
\end{frame}

\begin{frame}{Limites}
  \small
  \begin{itemize}
    \item[\textcolor{success}{\textbf{\checkmark}}] \textbf{Contrôle causal} :
      exécuté, \texttt{geo\_tol3\_permuted}. C'était la limite principale du
      rapport précédent.
    \item[\textcolor{failure}{\textbf{--}}] \textbf{Pas d'IC95} sur les écarts
      entre variantes : les deltas appariés sont calculés, le bootstrap groupé
      non. Plusieurs écarts sont du même ordre que la variance d'entraînement.
    \item[\textcolor{failure}{\textbf{--}}] \textbf{5 époques, pas 20.} Les
      variantes repartent d'un modèle convergé, mais l'optimum de validation de
      la baseline archivée est à l'époque 20.
    \item[\textcolor{failure}{\textbf{--}}] \textbf{\texttt{geo\_permuted} et
      \texttt{geo\_noise} non entraînés} : le $+0{,}0015$ de \texttt{geo} sur
      \texttt{cldice} n'est pas causalement attribué --- même si le contrôle du
      triplet \texttt{tol3} rend une conclusion différente très improbable.
    \item[\textcolor{failure}{\textbf{--}}] \textbf{\texttt{tol5} non exécuté} :
      la sensibilité au rayon de tolérance reste inconnue.
    \item[\textcolor{failure}{\textbf{--}}] \textbf{Évidence mono-échelle},
      accordée à une largeur moyenne peu représentative d'une partie du corpus.
    \item[\textcolor{failure}{\textbf{--}}] \textbf{Aucune évaluation
      multimodale}, ni sur ombres naturelles : l'argument central de la ligne
      n'est donc pas encore testé là où il pourrait tenir.
  \end{itemize}
\end{frame}

\begin{frame}{Suite}
  \small
  \begin{tabularx}{\linewidth}{c>{\raggedright\arraybackslash}X>{\raggedright\arraybackslash}p{3.35cm}}
    \toprule
    & Action & Justification \\
    \midrule
    \textbf{1} & \textbf{Adopter la perte tolérante par défaut.}
      \texttt{tol3} bat la baseline en strict \emph{et} en tolérant, pour un
      coût nul. & gain acquis, reproductible \\
    \textbf{2} & \textbf{Trancher la métrique du projet.} \texttt{clDice}
      domine à $k\ge1$ et en convention EUVIP ; \texttt{tol3} domine en
      strict. & les deux conventions ne classent pas pareil \\
    3 & \textbf{Porter l'expérience en multimodal sur FIND.} SAM 2 n'a pas
      accès à la portée : c'est le seul cadre où l'argument géométrique repose
      sur de l'information. & thèse de l'article ISPRS \\
    4 & Compléter les IC95 par bootstrap groupé, et le contrôle permuté de la
      famille clDice. & plusieurs écarts sont dans le bruit \\
    \bottomrule
  \end{tabularx}

  \vspace{0.24cm}
  \centering
  \statusbadge{failure}{à ne pas faire : augmenter la capacité de l'adapter, allonger l'entraînement, empiler un GNN}

  \vspace{0.18cm}
  {\small Sur Khánh Hà --- monomodal visible, annotations épaisses, baseline
  supervisée sur ce domaine même --- le Frangi généralisé n'apporte
  \textbf{aucune} information que la LoRA n'ait déjà apprise. Ce n'est pas un
  signal faible mal exploité : c'est l'absence de signal.}
\end{frame}

\appendix

\begin{frame}{Références}
  \small
  \begin{itemize}
    \item[\textbf{[1]}] Kirillov et al.,
      \href{https://arxiv.org/abs/2304.02643}{\emph{Segment Anything}},
      ICCV 2023.
    \item[\textbf{[2]}] Ravi et al.,
      \href{https://arxiv.org/abs/2408.00714}{\emph{SAM 2: Segment Anything in
      Images and Videos}}, 2024.
    \item[\textbf{[3]}] Ge et al.,
      \href{https://doi.org/10.1016/j.conbuildmat.2024.136573}{\emph{Fine-tuning
      vision foundation model for crack segmentation in civil
      infrastructures}}, Construction and Building Materials, 2024.
    \item[\textbf{[4]}] Shit et al.,
      \href{https://doi.org/10.1109/CVPR46437.2021.01629}{\emph{clDice --- a
      Novel Topology-Preserving Loss Function for Tubular Structure
      Segmentation}}, CVPR 2021.
    \item[\textbf{[5]}] Law et Chung, \emph{Three-dimensional curvilinear
      structure detection using optimally oriented flux}, ECCV 2008.
    \item[\textbf{[6]}] Kovesi, \emph{Symmetry and asymmetry from local phase},
      1997.
    \item[\textbf{[7]}] Frangi et al., \emph{Multiscale vessel enhancement
      filtering}, MICCAI 1998.
    \item[\textbf{[8]}] Hu et al.,
      \href{https://arxiv.org/abs/2106.09685}{\emph{LoRA: Low-Rank Adaptation
      of Large Language Models}}, 2022.
  \end{itemize}
\end{frame}

\end{document}
